توابع مقدماتی مختلط تعمیم توابع مقدماتی حقیقی‌اند (\lr{Elementary Complex Functions}).

\EMhold

\EMsection{تابع نمایی}{تابع نمایی}

\begin{EMunit}

تابع نمایی حقیقی \EMim{f(x)=e^x} دارای خواص زیر است:

\EMdisp{\begin{aligned}
&1)\quad f(x_1+x_2)=f(x_1)f(x_2)\qquad &&e^{x_1}e^{x_2}=e^{x_1+x_2}\\
&2)\quad \frac{d}{dx}f(x)=f(x)\qquad &&\frac{d}{dx}e^x=e^x\\
&3)\quad e^x=1+x+\frac{x^2}{2!}+\cdots+\frac{x^n}{n!}+\cdots
\end{aligned}}

\end{EMunit}

تابع نمایی مختلط را با الهام از خاصیت سوم تعریف می‌کنیم و سایر خواص را در مورد این تابع تحقیق می‌نمائیم.

\begin{EMdefinition}{تابع نمایی مختلط}

\EMdisp{w\EMbr =e^z\triangleq1+z+\frac{z^2}{2!}+\cdots+\frac{z^n}{n!}+\cdots}

\end{EMdefinition}

\begin{EMunit}

اگر در تعریف تابع نمایی مختلط به جای \EMim{z} مقدار \EMim{z=iy} قرار دهیم:

\EMdisp{\begin{aligned}
e^{iy}&=1+iy+\frac1{2!}(iy)^2+\frac1{3!}(iy)^3+\cdots\\
&=1+iy-\frac1{2!}y^2-\frac{i}{3!}y^3+\frac1{4!}y^4+\frac{i}{5!}y^5+\cdots\\
&=\underbrace{\left(1-\frac1{2!}y^2+\frac1{4!}y^4+\cdots\right)}_{\cos y}+i\underbrace{\left(y-\frac1{3!}y^3+\frac1{5!}y^5+\cdots\right)}_{\sin y}=\cos y+i\sin y
\end{aligned}}

\end{EMunit}

\begin{EMimportant}{}

\EMdisp{w\EMbr =e^z\EMbr =e^{x+iy}\EMbr =e^xe^{iy}\EMbr =e^x(\cos y+i\sin y)}

\end{EMimportant}

\begin{EMexample}{}

آیا تابع \EMim{w=e^z} تحلیلی است؟

\EMdisp{u(x,y)\EMbr =e^x\cos y,\qquad\EMbr  v(x,y)\EMbr =e^x\sin y}
\EMdisp{\frac{\partial u}{\partial x}\EMbr =e^x\cos y\EMbr =\frac{\partial v}{\partial y}\ \EMcheck \qquad\EMbr \qquad\EMbr  \frac{\partial u}{\partial y}\EMbr =-e^x\sin y\EMbr =-\frac{\partial v}{\partial x}\ \EMcheck }
چون شرایط کشی--ریمان در تمام صفحهٔ مختلط برقرار است، تابع \EMim{w=e^z} در تمام صفحهٔ مختلط مشتق‌پذیر است.

\end{EMexample}

\begin{EMexample}{}

مشتق تابع \EMim{w=e^z} را به‌دست آورید.

\EMdisp{\frac{d}{dz}e^z\EMbr =\frac{\partial u}{\partial x}+i\,\frac{\partial v}{\partial x}\EMbr =e^x\cos y+ie^x\sin y\EMbr =e^x(\cos y+i\sin y)\EMbr =e^z}

\end{EMexample}

\begin{EMimportant}{نکتهٔ مهم}

از این پس می‌توانیم هر عدد مختلط را به صورت نمائی نمایش دهیم:
\EMdisp{z\EMbr =x+iy\EMbr =|z|(\cos\theta+i\sin\theta)\EMbr =|z|\,e^{i\theta}}

\end{EMimportant}

\begin{EMunit}

مشاهده می‌شود که:
\EMdisp{e^{z+i2k\pi}\EMbr =e^ze^{i2k\pi}\EMbr =e^z\big[\cos(2k\pi)+i\sin(2k\pi)\big]\EMbr =e^z}
پس تابع \EMim{w=e^z} در صفحهٔ اعداد مختلط تابعی متناوب با دورهٔ تناوب \EMim{i2\pi} است.

\EMfigure{m10-exp-strip}{باند اصلی تابع نمایی: \EMim{-\pi<y\le\pi}}

\end{EMunit}

\EMhold

\EMsection{توابع مثلثاتی}{توابع مثلثاتی}

\begin{EMunit}

\EMdisp{e^{i\theta}\EMbr =\cos\theta+i\sin\theta\quad\EMbr \EMbr \Longrightarrow\quad\EMbr  e^{-i\theta}\EMbr =\cos\theta-i\sin\theta}
\EMdisp{\Longrightarrow\quad\EMbr \cos\theta\EMbr =\frac{e^{i\theta}+e^{-i\theta}}{2},\qquad\EMbr  \sin\theta\EMbr =\frac{e^{i\theta}-e^{-i\theta}}{2i}}

\end{EMunit}

\begin{EMdefinition}{کسینوس و سینوس مختلط}

\EMdisp{\cos z\triangleq\frac{e^{iz}+e^{-iz}}{2},\qquad\EMbr  \sin z\triangleq\frac{e^{iz}-e^{-iz}}{2i}}

\end{EMdefinition}

\begin{EMunit}

\EMdisp{\begin{aligned}
\cos z&=\frac12\Big[e^{i(x+iy)}+e^{-i(x+iy)}\Big]=\frac12\Big[e^{ix}e^{-y}+e^{-ix}e^{y}\Big]\\
&=\frac12\Big[e^{-y}(\cos x+i\sin x)+e^{y}(\cos x-i\sin x)\Big]\\
&=\cos x\left(\frac{e^y+e^{-y}}{2}\right)-i\sin x\left(\frac{e^y-e^{-y}}{2}\right)\\
&=\cos x\cosh y-i\sin x\sinh y
\end{aligned}}

\EMdisp{\begin{aligned}
\sin z&=\frac1{2i}\Big[e^{i(x+iy)}-e^{-i(x+iy)}\Big]=\frac1{2i}\Big[e^{ix}e^{-y}-e^{-ix}e^{y}\Big]\\
&=\frac1{2i}\Big[e^{-y}(\cos x+i\sin x)-e^{y}(\cos x-i\sin x)\Big]\\
&=\sin x\left(\frac{e^y+e^{-y}}{2}\right)+i\cos x\left(\frac{e^y-e^{-y}}{2}\right)\\
&=\sin x\cosh y+i\cos x\sinh y
\end{aligned}}

\end{EMunit}

\EMsubsection{تحلیلی بودن و مشتق‌پذیری \EMim{\cos z}}{تحلیلی بودن و مشتق‌پذیری cos z}

\EMdisp{\cos z\EMbr =u(x,y)+iv(x,y)\EMbr =\cos x\cosh y-i\sin x\sinh y,\qquad\EMbr  u\EMbr =\cos x\cosh y,\quad\EMbr  v\EMbr =-\sin x\sinh y}
\EMdisp{\frac{\partial u}{\partial x}\EMbr =-\sin x\cosh y\EMbr =\frac{\partial v}{\partial y}\ \EMcheck \qquad\EMbr \qquad\EMbr  \frac{\partial u}{\partial y}\EMbr =\cos x\sinh y\EMbr =-\frac{\partial v}{\partial x}\ \EMcheck }
پس این تابع در همهٔ صفحهٔ مختلط تحلیلی است. مشتق این تابع را توسط یکی از روش‌های ذکر شده به‌دست می‌آوریم:
\EMdisp{f'(z)\EMbr =\frac{d}{dz}\cos z\EMbr =\frac{\partial u}{\partial x}+i\,\frac{\partial v}{\partial x}\EMbr =-\sin x\cosh y-i\cos x\sinh y\EMbr =-\sin z}

\EMsubsection{تحلیلی بودن و مشتق‌پذیری \EMim{\sin z}}{تحلیلی بودن و مشتق‌پذیری sin z}

\EMdisp{\sin z\EMbr =u(x,y)+iv(x,y)\EMbr =\sin x\cosh y+i\cos x\sinh y,\qquad\EMbr  u\EMbr =\sin x\cosh y,\quad\EMbr  v\EMbr =\cos x\sinh y}
\EMdisp{\frac{\partial u}{\partial x}\EMbr =\cos x\cosh y\EMbr =\frac{\partial v}{\partial y}\ \EMcheck \qquad\EMbr \qquad\EMbr  \frac{\partial u}{\partial y}\EMbr =\sin x\sinh y\EMbr =-\frac{\partial v}{\partial x}\ \EMcheck }
پس این تابع در همهٔ صفحهٔ مختلط تحلیلی است و
\EMdisp{f'(z)\EMbr =\frac{d}{dz}\sin z\EMbr =\frac{\partial u}{\partial x}+i\,\frac{\partial v}{\partial x}\EMbr =\cos x\cosh y-i\sin x\sinh y\EMbr =\cos z}

\begin{EMimportant}{}

\EMdisp{(\sin z)'\EMbr =\cos z,\qquad\EMbr  (\cos z)'\EMbr =-\sin z}

\end{EMimportant}

\EMhold

\EMsubsection{بی‌کرانی \EMim{\sin z} و \EMim{\cos z}}{بی‌کرانی sin z و cos z}

\begin{EMunit}

\EMdisp{|\cos z|\EMbr =\sqrt{\cos^2x\cosh^2y+\sin^2x\sinh^2y}\EMbr =\sqrt{\cos^2x+\sinh^2y},\qquad\EMbr (\cosh^2y\EMbr =1+\sinh^2y,\ \ \sin^2x\EMbr =1-\cos^2x)}
\EMdisp{|\sin z|\EMbr =\sqrt{\sin^2x\cosh^2y+\cos^2x\sinh^2y}\EMbr =\sqrt{\sin^2x+\sinh^2y},\qquad\EMbr (\cosh^2y\EMbr =1+\sinh^2y,\ \ \cos^2x\EMbr =1-\sin^2x)}

\end{EMunit}

از آنجائی که \EMim{\sinh y} تابعی بی‌کران است، پس توابع \EMim{\cos z} و \EMim{\sin z} توابعی \textbf{بی‌کران} خواهند بود.

\EMhold

\EMsubsection{تناوب}{تناوب}

\begin{EMunit}

\EMdisp{\begin{aligned}
\cos(z+2\pi)&=\cos(x+iy+2\pi)=\cos(x+2\pi)\cosh y-i\sin(x+2\pi)\sinh y\\
&=\cos x\cosh y-i\sin x\sinh y=\cos z
\end{aligned}}
به همین ترتیب \EMim{\sin(z+2\pi)=\sin z}.

\EMfigure{m10-trig-strip}{باند اصلی توابع \EMim{\sin z} و \EMim{\cos z}: \EMim{-\pi<x\le\pi}}

\end{EMunit}

\EMhold

\EMsubsection{نکته}{نکته}

\begin{EMunit}

\EMdisp{\cos(iy)\EMbr =\cosh y,\qquad\EMbr  \sin(iy)\EMbr =i\sinh y}
\EMdisp{\cos(x+iy)\EMbr =\cos x\cos(iy)-\sin x\sin(iy)\EMbr =\cos x\cosh y-i\sin x\sinh y}
\EMdisp{\sin(x+iy)\EMbr =\sin x\cos(iy)+\cos x\sin(iy)\EMbr =\sin x\cosh y+i\cos x\sinh y}

\end{EMunit}

سایر توابع مثلثاتی:
\EMdisp{\tan z\triangleq\frac{\sin z}{\cos z},\qquad\EMbr  \cot z\triangleq\frac{\cos z}{\sin z},\qquad\EMbr  \sec z\triangleq\frac1{\cos z},\qquad\EMbr  \csc z\triangleq\frac1{\sin z}}
\EMdisp{(\tan z)'\EMbr =\sec^2z,\quad\EMbr \dots}

\EMhold

\EMsection{توابع هذلولی‌گون}{توابع هذلولی‌گون}

\begin{EMunit}

\EMdisp{\cosh z\triangleq\frac{e^z+e^{-z}}{2},\qquad\EMbr  \sinh z\triangleq\frac{e^z-e^{-z}}{2},\qquad\EMbr  e^z\EMbr =e^x(\cos y+i\sin y)}

\end{EMunit}

\begin{EMexercise}{}

تحقیق کنید:
\EMdisp{\cosh z\EMbr =\cosh x\cos y+i\sinh x\sin y,\qquad\EMbr  \sinh z\EMbr =\sinh x\cos y+i\cosh x\sin y}

\end{EMexercise}

\begin{EMexercise}{}

تحقیق کنید توابع \EMim{\cosh z} و \EMim{\sinh z} در تمام صفحهٔ مختلط تحلیلی‌اند و مشتق آنها از طریق روابط زیر به دست می‌آید:
\EMdisp{(\cosh z)'\EMbr =\sinh z,\qquad\EMbr  (\sinh z)'\EMbr =\cosh z}

\end{EMexercise}

\begin{EMunit}

\EMdisp{\tanh z\EMbr =\frac{\sinh z}{\cosh z},\quad\EMbr  \coth z\EMbr =\frac{\cosh z}{\sinh z},\quad\EMbr  \operatorname{sech}z\EMbr =\frac1{\cosh z},\quad\EMbr  \operatorname{csch}z\EMbr =\frac1{\sinh z}}

\EMdisp{\cosh(z+i2\pi)\EMbr =\cosh(z),\qquad\EMbr  \sinh(z+i2\pi)\EMbr =\sinh(z)}

\EMfigure{m10-hyp-strip}{باند اصلی توابع هذلولی‌گون: \EMim{-\pi<y\le\pi}}

\end{EMunit}

\EMhold

\EMsection{تابع لگاریتم}{تابع لگاریتم}

\begin{EMunit}

\EMdisp{w\EMbr =\ln z,\qquad\EMbr  z\EMbr =|z|e^{i\arg z}\EMbr =e^w\EMbr =e^{\ln|z|}e^{i\arg z}}
\EMdisp{w\EMbr =\ln z\EMbr =\ln|z|+i\arg z\EMbr =\ln\!\big(x^2+y^2\big)^{1/2}+i\tan^{-1}\frac yx}
\EMdisp{u(x,y)\EMbr =\frac12\ln\!\big(x^2+y^2\big),\qquad\EMbr  v(x,y)\EMbr =\tan^{-1}\frac yx}
\EMdisp{\frac{\partial u}{\partial x}\EMbr =\frac{x}{x^2+y^2}\EMbr =\frac{\partial v}{\partial y}\ \EMcheck \qquad\EMbr \qquad\EMbr  \frac{\partial u}{\partial y}\EMbr =\frac{y}{x^2+y^2}\EMbr =-\frac{\partial v}{\partial x}\ \EMcheck }

\end{EMunit}

\begin{EMunit}

بنابراین تابع لگاریتم در تمام صفحهٔ مختلط بجز مبدأ تحلیلی است.

\EMdisp{\frac{d}{dz}\ln z\EMbr =\frac{\partial u}{\partial x}+i\,\frac{\partial v}{\partial x}\EMbr =\frac{x}{x^2+y^2}+i\,\frac{-y}{x^2+y^2}\EMbr =\frac{\bar z}{z\bar z}\EMbr =\frac1z}

\EMdisp{z\EMbr =re^{i\theta}\EMbr =re^{i(\theta+2k\pi)}}

\end{EMunit}

\begin{EMunit}

برای آنکه تعریف تابع نقض نشود برد این تابع ناحیهٔ نشان داده شده در شکل است.

\EMfigure{m10-log-range}{برد تابع لگاریتم: نوار \EMim{-\pi<\operatorname{Im}w\le\pi}}

\end{EMunit}

خواص زیر برای تابع لگاریتم برقرارند:
\EMdisp{\ln(z_1\times z_2)\EMbr =\ln z_1+\ln z_2,\qquad\EMbr  \ln\!\left(\frac{z_1}{z_2}\right)\EMbr =\ln z_1-\ln z_2,\qquad\EMbr  \ln(z^c)\EMbr =c\ln z}
